\chapter{The Two Roles of Classical Mechanics}
\label{tr}

\section{Introduction}

Classical mechanics occupies a peculiar position among physical theories. It is,
first, a complete and self-contained subject: a body of concepts, equations, and
results describing the motion of bodies under forces, valid and useful on its own
terms, with no reference whatsoever to the Planck constant $\hbar$ or the speed of
light $c$. Celestial mechanics, rigid-body dynamics, the theory of small
oscillations, continuum mechanics, and the modern theory of dynamical systems and
chaos are all classical mechanics pursued for its own sake, and none of them is
waiting to be corrected by a more fundamental theory in order to be meaningful.

Second, and for the purposes of this book more importantly, classical mechanics is
the \emph{base} on which relativistic and quantum mechanics are built. Special and
general relativity are commonly introduced as corrections to Newtonian mechanics
that become significant when velocities approach $c$ or fields become strong;
quantum mechanics is commonly introduced as a modification that becomes significant
at the scale of $\hbar$. In both cases, the classical theory is treated as a known,
well-understood limit --- the solid ground from which the generalization departs
and to which it must reduce.

It is tempting to think that once this second, subordinate role is granted, nothing
further needs to be said about classical mechanics: its job is simply to sit still
and be a limit. The purpose of this chapter, and of the chapters immediately
following it, is to argue that this is false in two related ways.

\begin{enumerate}
  \item[(a)] Even restricted to the narrow task of serving as the base for
  relativistic and quantum generalization, classical mechanics is not a settled,
  transparent starting point. Producing the structure that gets generalized ---
  a well-defined phase space, a Hamiltonian, a consistent variational principle
  --- is itself a nontrivial achievement that fails, or becomes ambiguous, for
  entire classes of ordinary classical systems.
  \item[(b)] Several phenomena that are routinely presented as characteristically
  relativistic or characteristically quantum-mechanical --- an underdetermination
  of physical ontology by the observable facts, an ambiguity in the passage from
  classical description to quantized theory, a geometric ``phase'' with no
  classical counterpart, an apparent conflict between microscopic reversibility
  and macroscopic causal, dissipative behavior --- already occur inside classical
  mechanics, before relativity or $\hbar$ ever enter the discussion. When such a
  paradox survives the excision of everything relativistic or quantum, its source
  cannot honestly be attributed to relativity or to quantum theory.
\end{enumerate}

Section~\ref{tr:st} says a brief word about the first, standalone role
of classical mechanics before setting it aside. Section~\ref{tr:ba} examines
what is actually required of classical mechanics for it to serve as a base for
generalization, and where that requirement is met only imperfectly.
Section~\ref{tr:pa} works through four examples --- Hertz's forceless
mechanics, nonholonomic constraints, shape-changing bodies, and open dissipative
systems --- each of which is treated at length, in its own right, later in this
book; here they are introduced specifically to show that a paradox usually filed
under ``quantum'' or ``relativistic'' foundations already has a classical
instance.

\section{Classical Mechanics as a Subject in Its Own Right}
\label{tr:st}

It is worth pausing on the first role only long enough to insist that it is real.
The three-body problem, the precession of orbits, the stability of the solar
system, the dynamics of rigid bodies and gyroscopes, the propagation of waves in
elastic media, the onset of turbulence, and the sensitive dependence on initial
conditions studied by dynamical-systems theory are all substantial, still
partially open, areas of physics conducted entirely within the classical
framework. None of this material is a mere warm-up exercise for quantum mechanics,
and a physicist could spend a full career inside it without ever needing $\hbar$.
This book is not about that role. It is mentioned here only to make clear that
when the rest of this book treats classical mechanics as a base to be generalized,
it is not thereby claiming that this exhausts what classical mechanics is.

\section{Classical Mechanics as the Base of Generalization}
\label{tr:ba}

The textbook picture of generalization is familiar. Quantization is presented as
a recipe applied to a classical Hamiltonian system: promote canonical coordinates
$(q,p)$ to operators $(\hat q,\hat p)$, and promote the Poisson bracket to the
commutator,
\begin{equation}
  \{f,g\} \;\longmapsto\; \frac{1}{i\hbar}\,[\hat f,\hat g] \, .
\end{equation}
The inverse limit, $\hbar \to 0$, recovers classical mechanics via the WKB
approximation and Ehrenfest's theorem. Similarly, special relativity is presented
as reducing to Newtonian mechanics for $v \ll c$, and general relativity as
reducing to Newtonian gravitation in the weak-field, slow-motion limit.

Both halves of this picture presuppose that classical mechanics has already
delivered, for the system in question, a specific and unambiguous piece of
structure: a phase space with a well-defined symplectic (or Poisson) structure,
a Hamiltonian generating the correct equations of motion, and a family of
observables closed under the bracket. Only given that structure does
``quantize it'' or ``take the classical limit'' pick out a definite theory.

For a mass point moving under a potential, or for a system subject only to
holonomic constraints, this presupposition is satisfied and the textbook picture
is accurate. But it is easy to overstate how general this good case is. As soon
as one leaves it --- constraints that cannot be integrated to relations among
coordinates alone, systems that exchange energy or momentum with an unmodeled
environment, systems whose configuration space itself is not fixed once and for
all --- the classical structure needed as input to generalization becomes
ambiguous, multiple, or simply unavailable in closed form. In those cases the
question ``what is the classical mechanics of this system'', asked with no
reference to $\hbar$ or $c$ at all, already has more than one defensible answer.
That is the sense in which classical mechanics, even in its purely subordinate
role as a base, is \emph{far from trivial}: the base itself has to be constructed,
and the construction is not always unique.

\section{Paradoxes with Classical Roots}
\label{tr:pa}

The four cases below are developed fully in later chapters of this book. They
are introduced here, briefly and together, because their common moral is best
seen by contrast: in each case, a feature usually advertised as a signature of
quantum theory or of relativity turns out to have an exact, or near-exact,
analogue that requires nothing beyond Newton, Lagrange, and Hamilton.

\subsection{Hertz's forceless mechanics and the underdetermination of ontology}

In his \emph{Principles of Mechanics}, Heinrich Hertz reformulated dynamics
without the concept of force at all. Observable bodies are treated as coupled to
unobservable ``hidden masses'' through rigid, holonomic constraints, and the
combined system moves along a geodesic of a metric built from the mass
distribution on an enlarged configuration space. Suitably arranged, Hertz's
constrained geodesic motion reproduces the trajectories that Newtonian mechanics
attributes to the action of a force. The two formulations are empirically
equivalent: no experiment performed on the observable bodies alone can decide
between ``a force is acting'' and ``a constraint is enforcing geodesic motion in
a larger space.''

This is a case of underdetermination of theory --- indeed of ontology --- by
observation, arising entirely within classical mechanics. It is normally quantum
mechanics that is credited with forcing physicists to confront multiple,
empirically indistinguishable accounts of what is ``really'' happening beneath
the observed statistics: Bohmian trajectories, Everettian branches, and
epistemic and objective-collapse readings of the wavefunction are all consistent
with the same experimental record. Hertz shows that the same structural
situation --- one set of observations, several inequivalent stories about the
underlying ontology --- already existed in nineteenth-century particle mechanics,
with no probability amplitude in sight.

\subsection{Nonholonomic constraints and the ill-posedness of quantization}

A constraint of the Pfaffian form
\begin{equation}
  \sum_i a_i(q)\,\dot q_i = 0
\end{equation}
is \emph{nonholonomic} when it is not integrable to a relation $f(q)=0$ among the
coordinates alone --- the standard example being a disk or coin rolling without
slipping on a plane. Such systems present a genuine fork in the road. Applying
d'Alembert's principle directly, with the constraint enforced at the level of
virtual displacements, yields the correct equations of motion, confirmed by
experiment. Applying Hamilton's variational principle instead --- extremizing the
action with the constraint imposed by a Lagrange multiplier, as one legitimately
may for holonomic constraints --- yields a different, and physically incorrect,
set of equations. This is the long-standing nonholonomic-versus-vakonomic
controversy: two superficially reasonable derivations from the same starting
Lagrangian and the same constraint disagree about the dynamics, and only one of
them is right.

One consequence is that nonholonomic systems generically fail to admit a
straightforward canonical Hamiltonian formulation on the naive phase space:
the constrained bracket structure one is forced to use is not the canonical
Poisson bracket, and which reduced structure is the correct one to use is itself
a nontrivial derivation, system by system. This matters directly for the
generalization picture of Section~\ref{tr:ba}: canonical quantization
needs a Poisson bracket to promote to a commutator, and here the classical
bracket is already ambiguous before any promotion is attempted. The
operator-ordering problem --- the fact that a classical observable such as $qp$
does not determine a unique quantum operator, since $\hat q\hat p$, $\hat p\hat
q$, and $\tfrac12(\hat q\hat p+\hat p\hat q)$ all reduce to it as $\hbar\to0$
--- is usually presented as an intrinsically quantum indeterminacy. In the
nonholonomic case the indeterminacy has already begun classically, in the
disagreement over what the correct pre-quantization bracket even is.

\subsection{Shape-changing bodies and classical geometric phase}

A body that changes its own shape --- the paradigm case is a cat, initially
oriented feet-up, that reorients itself in mid-air with zero net angular
momentum by cyclically changing its shape --- cannot be treated by rigid-body
mechanics. One instead works on a \emph{shape space} $\mathcal S$ (the space of
internal configurations, with rigid motions quotiented out) and finds that the
conservation of angular momentum, together with the equations of motion,
defines a connection on a principal bundle over $\mathcal S$ with structure
group $SO(3)$. A closed loop in shape space --- the cat returns to its initial
shape --- need not be a closed loop in the full configuration space: the body
can acquire a net rotation purely as a holonomy of this connection, with no
external torque and no change in angular momentum at any instant. The same
mechanism governs low-Reynolds-number swimmers and, more generally, any system
in which a periodic change of internal, ``slow'' variables produces a net shift
in an external, cyclic variable; the resulting angle is called the Hannay angle.

The geometric or Berry phase acquired by a quantum system whose Hamiltonian is
slowly and cyclically varied is frequently introduced as a paradigmatically
quantum-mechanical effect, precisely because it has no counterpart in the naive
classical picture of a particle following a trajectory. The falling cat and the
Hannay angle show that this is too strong a claim about geometric phases in
general: the underlying mechanism --- holonomy of a connection around a loop in
parameter or shape space, decoupled from the local dynamics along the way --- is
already fully present, and was fully worked out, in ordinary classical
mechanics. What is genuinely new in the quantum case is the specific bundle and
connection defined by the adiabatic theorem applied to a Hilbert space of
states; the geometric phenomenon of holonomy itself is not new.

\subsection{Open systems, radiation reaction, and the arrow of time}

Two classical difficulties are worth setting side by side. The first is
statistical: Boltzmann's $H$-theorem derives macroscopic irreversibility from
microscopically time-reversible mechanics, and was met immediately by
Loschmidt's reversibility objection (reversing all velocities should reverse
the approach to equilibrium) and Zermelo's recurrence objection (Poincar\'e
recurrence guarantees a return arbitrarily close to any earlier state). Both
objections are correct as stated, and reconciling them with the observed arrow
of time is a genuinely hard, and still actively discussed, problem --- entirely
classical, with no reference to measurement or wavefunction collapse.

The second is dynamical: an accelerating point charge radiates, and Newton's
second law must be supplemented by the Abraham--Lorentz self-force,
\begin{equation}
  m\dot v = F_{\text{ext}} + \frac{2}{3}\frac{e^2}{c^3}\,\dot a \, ,
\end{equation}
a third-order correction proportional to the derivative of the acceleration.
This equation admits runaway solutions, in which the particle self-accelerates
exponentially with no external force at all; demanding that solutions remain
bounded, instead, forces the particle to begin accelerating \emph{before} the
external force is applied --- a pre-acceleration that appears to violate
causality, entirely within classical electrodynamics.

Both irreversibility and an apparent tension with causal, deterministic
evolution are frequently treated as problems that quantum theory introduces
--- through the measurement postulate, decoherence, or the collapse of the
wavefunction --- or that relativity introduces, through its reworking of
simultaneity. Loschmidt, Zermelo, and Abraham--Lorentz show that a purely
classical system, governed by Newton's laws and Maxwell's equations and nothing
else, already exhibits dissipative, apparently irreversible, and apparently
acausal behavior that has never been fully tamed.

\section{Looking Ahead}

Each of the four cases above is the subject of a full chapter later in this
book, where the mathematics is developed in detail and the historical debate is
traced more carefully. The point of gathering them here, in outline, is
methodological: before a ``paradox'' encountered in relativity or in quantum
mechanics is attributed to relativity or to quantum mechanics as such, it is
worth checking whether classical mechanics, unaided, already contains a version
of it. Where it does, what is new in the generalization is not the paradox but,
at most, its dressing --- and the genuinely new content of the generalization
has to be located elsewhere.

